% !TEX root = main.tex
% -----------------------------------------------------------------------------
% Chapter 2: Literature Review
% -----------------------------------------------------------------------------

\chapter{Literature Review}
\label{ch:literature-review}

\section{Overview of Missing Data in Sensor Time Series}

Missing data is a familiar problem in statistics, but it takes on a slightly different character in sensor networks. A missing value in a questionnaire may mean that a participant skipped a question. A missing value in an IoT temperature stream may mean that the wireless link dropped for ten minutes, the battery dipped below a usable level, the device restarted, the timestamp parser failed, or the sensor briefly produced readings that were so unrealistic that they had to be treated as unusable. The data point is gone either way, but the practical story behind the gap is different. That story matters because it can change what kind of imputation is defensible.

The literature on IoT and environmental sensing describes missing data as a recurring data-quality problem rather than a rare accident. Surveys of sensor imputation report that missingness can arise from hardware faults, communication problems, synchronization errors, energy constraints, calibration drift, and data-management failures \cite{Adhikari2022ACS,Agbo2022MissingDI,Osman2018ASO}. In water distribution, air-quality monitoring, smart grids, and low-cost environmental sensing, incomplete streams are not just inconvenient; they can weaken the reliability of forecasting, calibration, anomaly detection, and decision support \cite{Osman2018ASO,Schreiber2023DataIT,Agbo2022MissingDI}. The same concern motivates this thesis. A temperature series may look simple at first glance, but once values disappear in blocks or at the edge of the sequence, the reconstruction problem becomes less forgiving.

\subsection{Sources and patterns of missingness}

Sensor missingness is shaped by both the device and the environment around it. A sensor node may stop reporting because of power loss, wireless interference, gateway failure, physical damage, memory constraints, maintenance downtime, or temporary network congestion. In outdoor or semi-controlled deployments, weather and location also matter: heavy rain, heat, enclosure issues, or poor placement may affect the device exactly when the measurement is most useful. In industrial and smart-grid settings, missingness may also appear when systems are reconfigured, when meters are replaced, or when measurements are rejected by validation rules \cite{Schreiber2023DataIT,Wu2022DataIF}.

The pattern of the missing values is at least as important as the proportion missing. A dataset with 30\% missing values scattered one point at a time is very different from one with a single 30\% missing block. Short random gaps often leave usable information on both sides. Long gaps remove local context. Boundary gaps are stranger still: when the first part of a sequence is missing, there is no past context; when the final part is missing, there is no future context. Studies on long-gap IoT and multivariate sensor imputation point out that larger contiguous gaps often require methods that can exploit wider temporal, spatial, or seasonal information rather than only nearest-neighbour timestamps \cite{Ahmed2022LongGM,Wu2022DataIF,Liu2020}.

This distinction explains why many recent evaluations move beyond one generic random missingness setting. Time-series benchmarking work has shown that the same method can look strong under one missingness pattern and ordinary under another \cite{Toye2025BenchmarkingMD}. That finding may sound obvious, but it is easy to forget in practice. If a paper reports excellent performance after deleting scattered observations at random, it has not necessarily shown that the method can recover the start of a recording, a long communication outage, or a full missing afternoon in a field sensor deployment.

\subsection{Missing-data mechanisms}

Classical missing-data theory usually separates missingness into three mechanisms: missing completely at random (MCAR), missing at random (MAR), and missing not at random (MNAR). In MCAR, the probability that a value is missing is unrelated to the observed or unobserved data. In MAR, the probability of missingness may depend on observed variables. In MNAR, the missingness depends on the missing value itself or on unobserved information \cite{Little2019,Jamshidian2014,Pham2022}. These categories are useful because they remind researchers that missingness is not only a numerical inconvenience. It is also an assumption about how the data became incomplete.

Sensor data often sits awkwardly between these neat categories. A random packet loss event can resemble MCAR if it is unrelated to the temperature being measured. A battery issue may look closer to MAR if voltage readings or device status are observed and can explain the missingness. A sensor that fails during extreme heat, or drops out when humidity enters the enclosure, starts to resemble MNAR because the missingness may be linked to the very condition the analyst wants to measure. In many real deployments, the mechanism is only partly known. The analyst may see the timestamps and values, but not the exact reason a device stopped sending.

That uncertainty makes synthetic masking both useful and limited. It is useful because it creates known ground truth: complete observations can be temporarily hidden, imputed, and then compared against their true values. It is limited because synthetic masks rarely capture every real failure mode. Farjallah et al. note that evaluation without ground truth becomes difficult in genuine missing-data settings, while benchmark studies often rely on artificial removal from complete data \cite{Farjallah2025EvaluationOM,Toye2025BenchmarkingMD}. The present thesis follows the controlled-masking route because it allows fair comparison, but the literature suggests a cautious interpretation: good performance under artificial masks may suggest practical value, not prove it for every real outage.

\subsection{Implications for downstream analysis}

Missing sensor values can distort later analysis in several ways. Some algorithms reject missing inputs outright. Others silently handle them in ways that may not match the analyst's intention. For time series, deletion is especially unattractive because removing timestamps changes the spacing of the sequence. A forecasting model trained after listwise deletion may no longer see the actual temporal rhythm of the system. A calibration model may learn from a subset that over-represents calm periods and under-represents unusual ones. An anomaly detector may miss exactly the kind of deviation it is supposed to detect.

The effect of imputation also depends on what the completed data will be used for. If the goal is to estimate a daily average temperature, a simple method may be adequate. If the goal is to detect a short heating fault or recover a sharp overnight change, the same method may smooth away the feature of interest. Reviews repeatedly warn that imputation should be selected with the data type, missingness pattern, downstream task, and evaluation objective in mind \cite{Alwateer2024MissingDI,Adnan2022ARO,Miao2023AnES}. This is a modest claim, but an important one: an imputer is not just a preprocessing button. It becomes part of the modelled signal.

In IoT settings, computation is another downstream issue. Edge devices and gateways may have limited memory and processing time, so a method that is attractive in an offline notebook may be harder to justify in a deployed pipeline. Erhan et al. and related IoT-edge work show the appeal of methods that can clean data close to the source, but they also make the trade-off visible: accuracy, memory use, latency, and implementation complexity all matter \cite{Erhan2021CriticalCO,Agbo2020BestFM}. For this thesis, runtime is not the main scientific contribution, yet it is still worth measuring because a hybrid LSTM-VAE carries more machinery than mean imputation or interpolation.

\begin{table}[ht]
	\centering
	\small
	\begin{tabular}{p{0.23\linewidth}p{0.34\linewidth}p{0.33\linewidth}}
		\toprule
		\textbf{Literature strand} & \textbf{Typical methods} & \textbf{Relevance to this thesis} \\
		\midrule
		Statistical imputation & Mean, median, mode, hot-deck, EM, MICE, interpolation & Provides transparent baselines and highlights what can be achieved without a neural model \cite{Alwateer2024MissingDI,VanBuuren2011,Dempster1997}. \\
		Machine-learning imputation & KNN, regression, random forest, gradient boosting, iterative regression & Tests whether engineered temporal features and non-neural learners already capture enough structure \cite{Raj2023MachineLB,Miao2023AnES,Lee2025ComparisonOM}. \\
		Deep sequence imputation & LSTM, Bi-LSTM, recurrent autoencoders, attention-based models & Motivates sequence models for long and structured gaps, especially where local interpolation is weak \cite{Hochreiter1997,Aqilah2025ASR,Lyu2022}. \\
		Generative imputation & VAE, denoising autoencoder, GAN, diffusion model & Supports latent reconstruction, uncertainty-aware modelling, and residual refinement, but raises training and evaluation issues \cite{Okafor2021,Ou2024MissingDataIW,Yoon2018GAINMD,Zhang2024DiffPuterED}. \\
		\bottomrule
	\end{tabular}
	\caption{Main strands of missing-data imputation literature used to position the proposed hybrid LSTM-VAE approach.}
	\label{tab:lit-strands}
\end{table}

\section{Classical Statistical Imputation Methods}

Classical imputation methods remain widely used because they are simple, quick, and easy to explain. This is sometimes treated as a weakness, but it can be a strength in applied work. A field engineer, supervisor, or domain reviewer can usually understand mean replacement or linear interpolation without needing to inspect training curves or latent variables. The real question is not whether these methods are simple. It is whether their assumptions are good enough for the missingness pattern at hand.

\subsection{Mean, median, and mode-based approaches}

Mean, median, and mode imputation replace missing values with a single representative statistic calculated from the observed data. Mean imputation uses the arithmetic average, median imputation uses the middle observed value, and mode imputation uses the most frequent observed value. These methods are still common in comparative studies because they provide a minimal baseline: if a proposed model cannot outperform central-tendency replacement, something is probably wrong with either the model or the evaluation design.

For sensor time series, however, constant replacement is usually a blunt instrument. A temperature sequence has temporal shape: morning readings may differ from afternoon readings, and seasonal or daily patterns can make the same numerical value plausible at one time and implausible at another. Filling every gap with the same mean value flattens that structure. Median imputation can be less sensitive to outliers, which is useful when sensor spikes exist, but it still ignores ordering. Mode imputation is often even less natural for continuous temperature readings because the most frequent value may reflect rounding rather than physical centrality.

The literature generally treats these methods as convenient but limited. Reviews of missing-data methods note that simple imputation can reduce variance, weaken relationships among variables, and create overconfident completed datasets \cite{Little2019,Alwateer2024MissingDI}. In time-series sensor data, the problem is visible in a more concrete way: a flat segment appears where a curve should be. That flatness may not harm a coarse summary, but it can damage local reconstruction. This is why constant methods are useful in this thesis mainly as reference points, not as serious candidates for the best reconstruction.

\subsection{Forward fill, backward fill, and interpolation}

Time-aware statistical methods take one step closer to the nature of sensor data. Forward fill, often called last observation carried forward (LOCF), replaces a missing value with the most recent observed value. Backward fill or last observation carried backward (LOCB) does the opposite, using the next observed value. Linear interpolation estimates values along a straight line between two observed endpoints. These methods still make simple assumptions, but at least they use the sequence order.

Forward fill can work surprisingly well when a sensor changes slowly and gaps are short. A temperature reading one minute from now may be close to the current reading, especially indoors or in a stable outdoor period. But if the gap spans a transition, forward fill can create a staircase pattern. Backward fill has the same issue in reverse. Interpolation is often more visually plausible because it connects the two sides of a gap, and for smooth temperature curves it can be a difficult baseline to beat. Several time-series imputation comparisons show that interpolation and related forecasting methods remain competitive for many real-world sensor datasets \cite{Ahn2022ComparisonOM,RoseroMontalvo2023TimeSeriesFT,Toye2025BenchmarkingMD}.

The trouble begins when interpolation is asked to do too much. A straight line across a long missing block may erase a local peak, miss a sudden drop, or invent a smooth transition that never happened. At the start or end of a sequence, ordinary interpolation also loses one of its anchors. Boundary handling then becomes a secondary imputation decision: should the method carry the nearest value, use a global statistic, fit a trend, or search for similar seasonal points? The answer is not automatic. For this reason, the proposed model in this thesis uses interpolation mainly for MIDDLE and RANDOM masks, where neighbouring context is usually available, and uses a seasonal KNN prior for START and END masks, where interpolation is less naturally supported.

\subsection{Strengths and limitations of simple methods}

The main strengths of simple imputation methods are transparency, speed, and reproducibility. They require little tuning. They are easy to implement. They also expose a useful practical fact: many sensor signals are smooth enough that a simple method can be quite strong. It would be a mistake to dismiss interpolation only because it is not a deep model. In fact, a neural model that cannot beat linear interpolation on a smooth internal gap may not be adding much value.

The limitations are equally clear. Simple methods tend to ignore wider context, uncertainty, nonlinear temporal structure, and repeated seasonal patterns. Constant methods damage the local shape of the series. Carry-forward methods can freeze the signal. Linear interpolation assumes a straight path between observed endpoints. These assumptions may be acceptable for short gaps, but they become less plausible as gap length increases or when missingness occurs near sequence boundaries. Studies of large-gap sensor imputation and IoT missingness repeatedly point to this problem \cite{Wu2022DataIF,Ahmed2022LongGM,Liu2020}.

There is also a subtle evaluation issue. If researchers compare a deep model only against mean, median, or mode imputation, the deep model may look better than it really is. A fairer comparison includes time-aware baselines such as interpolation and forward/backward fill, because those are the methods a practitioner would likely try first. This thesis follows that view by including both weak lower-bound baselines and stronger time-aware statistical baselines.

\section{Machine Learning-Based Imputation Methods}

Machine-learning imputers sit between simple statistical rules and deep sequence models. They can learn relationships from observed data, but they usually need the time series to be converted into features. In univariate sensor imputation, this often means building lagged inputs: previous readings, rolling statistics, calendar features, mask indicators, or local trend features. The model then predicts the missing value from those features.

\subsection{K-nearest neighbours and distance-based methods}

K-nearest neighbours (KNN) is a natural starting point for machine-learning imputation because it is easy to understand. It imputes a missing value by finding similar observed cases and averaging or otherwise combining their values. In tabular data, similarity may be based on other variables. In time series, it may be based on lagged readings, time-of-day features, day-of-year features, or nearby sensors. KNN appears often in missing-data comparisons because it can capture local structure without fitting a complicated parametric model \cite{Raj2023MachineLB,Miao2023AnES}.

For temperature data, KNN can be interpreted in a fairly concrete way. If a reading is missing at 6:00 a.m., the method may look for previous windows that had similar recent temperature patterns or similar calendar positions. This resembles how a person might reason: "What happened on other mornings that looked like this?" That analogy is imperfect, but it explains why distance-based methods can be useful for seasonal or repeated signals.

The weakness is that similarity depends heavily on the chosen representation. If the lag window is too short, the model may miss slower patterns. If the feature scale is poor, one feature can dominate the distance. If gaps are long, many lagged values may themselves be imputed, so the neighbour search rests on partly artificial inputs. KNN can also become expensive for large datasets because prediction requires searching through stored examples. In edge or near-real-time settings, that cost may matter \cite{Erhan2021CriticalCO}.

\subsection{Regression-based and iterative imputers}

Regression-based imputation estimates missing values by learning a mapping from observed features to the target variable. Linear regression and ridge regression provide simple versions of this idea. They are attractive when the relation between lagged context and the missing value is mostly smooth or linear. Ridge regression, in particular, can handle correlated lag features more gracefully than ordinary least squares because the penalty shrinks unstable coefficients.

Iterative imputers extend this idea to multiple incomplete variables or lagged matrices. Multiple Imputation by Chained Equations (MICE) fills each incomplete variable using a model based on the others, cycling through variables repeatedly until the imputations settle \cite{VanBuuren2011}. Expectation-maximization (EM), although not always used as a direct single-value imputer in the same way, is another foundational approach for inference with incomplete data \cite{Dempster1997}. These methods come from a statistical tradition that pays close attention to assumptions and uncertainty. That tradition is worth keeping in view, even when the implementation in this thesis uses a simpler iterative Bayesian-ridge baseline.

For univariate time series, iterative imputation can be adapted by embedding the sequence into a lagged matrix. Each row represents a local window, and each column represents a lag position. The imputer then learns relationships among lag columns. This can work well when the series has repeated local structure. It may struggle, however, when the missing block is so large that many lagged columns are simultaneously incomplete. In that case, the imputer may recycle its own guesses.

\subsection{Tree-based and ensemble approaches}

Tree-based methods, including random forests, gradient boosting, and related ensemble models, are widely used for imputation because they can model nonlinear relationships and feature interactions without requiring the analyst to specify them directly. Studies of missing-data imputation often include random forest variants, gradient boosting, or missForest-style methods, and results commonly show that these models are competitive on structured tabular and sensor datasets \cite{Miao2023AnES,Lee2025ComparisonOM,Luo2022}.

In sensor settings, tree-based models can combine several cues: recent lags, rolling averages, time-of-day, day-of-year, mask indicators, and sometimes correlated sensor channels. This can be helpful when the missing value depends on a mixture of local trend and seasonal context. A histogram gradient boosting model, for example, may learn that a morning temperature rise behaves differently from an evening temperature drop, provided those patterns are represented in the features.

Their limitations are more practical than philosophical. Tree ensembles can be sensitive to feature design, missingness simulation, and train-test separation. They also produce point estimates unless additional uncertainty machinery is added. When the task is sequence reconstruction, a tree model predicts individual timestamps rather than generating a coherent curve unless the feature construction encourages temporal consistency. That does not make the approach weak; it simply means that a good error score may still need visual inspection.

\subsection{Comparative performance on time-series data}

Comparative studies do not point to one universally best imputation method. That is probably the most consistent finding in the literature. Miao et al. show broad variation across algorithms and datasets in an experimental survey \cite{Miao2023AnES}. Toye et al. also emphasize that time-series benchmarks should use realistic test cases because method rankings can shift with the dataset and missingness pattern \cite{Toye2025BenchmarkingMD}. In IoT and environmental sensing, Agbo et al. and Erhan et al. similarly show that performance depends on missingness rate, gap structure, and computational constraints \cite{Agbo2022MissingDI,Erhan2021CriticalCO}.

This has a direct consequence for the present study. A hybrid LSTM-VAE should not be evaluated only against weak statistical baselines. KNN regression, ridge regression, gradient boosting, and iterative imputation represent more serious alternatives. If the proposed model improves over those methods under the same masks and metrics, the claim is stronger. If it does not, that is still informative: it may suggest that the signal is sufficiently captured by lagged features or deterministic priors.

\section{Deep Learning-Based Imputation Methods}

Deep learning entered the missing-data literature partly because many real datasets have nonlinear structure that simple imputers do not capture well. Sequence models can learn temporal dependencies, autoencoders can learn compressed representations, and generative models can sample or reconstruct plausible missing values. Even so, the literature is less tidy than the headline suggests. Deep models can perform very well, but they can also be data-hungry, difficult to tune, computationally expensive, and surprisingly easy to over-praise when the baselines are weak.

\subsection{Recurrent neural networks for sequence reconstruction}

Recurrent neural networks are designed for ordered data, and LSTM networks were introduced to address the difficulty of learning long-range dependencies in ordinary recurrent networks \cite{Hochreiter1997}. In imputation, this matters because a missing value is often related not only to the immediately adjacent value but also to a wider temporal pattern. A temperature reading may reflect a daily cycle, a slow weather change, or a delayed response to heating and cooling. LSTM models can, in principle, learn some of these patterns from sequences.

Recent reviews of recurrent neural networks for imputation report growing use of RNN, LSTM, GRU, bidirectional, and attention-enhanced architectures \cite{Aqilah2025ASR}. Applied sensor studies also show why this family is attractive. Rosero-Montalvo et al. compare forecasting methods for IoT missing-data filling and report that different sensor channels may favour different methods, including LSTM for some signals \cite{RoseroMontalvo2023TimeSeriesFT}. Lyu et al. combine Bi-LSTM, attention, and transfer learning for large continuous missing blocks in water-quality data \cite{Lyu2022}. Vedavalli and Deepak use a spatial-temporal hierarchical LSTM approach for missing and erroneous IoT node data \cite{Vedavalli2023}.

There is a catch, though. RNN-based imputation can be sensitive to how the missing values are represented during training. Replacing missing points with zero may introduce artificial jumps. Training only on observed points may leave the model uncertain about long masked intervals. Bidirectional models can use both past and future context, but that is not always available in online settings. LSTM models also do not automatically know the calendar structure of a sensor series unless those features are supplied or learned indirectly. For the present thesis, these issues motivate standardisation, mask inputs, calendar features, and deterministic priors.

\subsection{Autoencoders and variational autoencoders}

Autoencoders learn to compress data into a latent representation and then reconstruct it. Denoising autoencoders adapt this idea by training the model to recover clean data from corrupted inputs. This is a close conceptual match to imputation: the observed sequence is incomplete or corrupted, and the model learns to reconstruct the missing parts. Position-encoding denoising autoencoders have been proposed for industrial process data, suggesting that time or position information can help the model distinguish one part of a sequence from another \cite{Ou2024MissingDataIW}.

Variational Autoencoders (VAEs) add a probabilistic latent representation. Rather than mapping an input to a single latent code, a VAE maps it to a latent distribution and regularises that distribution toward a prior. In imputation, this can be useful because missing values are uncertain. The model is not merely filling a blank cell; it is producing a plausible reconstruction conditioned on incomplete evidence. VAE-based imputation has been applied to IoT sensor calibration, where Okafor and Delaney compare VAE with KNN, MICE, random forest imputation, and other methods \cite{Okafor2021}. More recent work has also explored deep VAEs for process data quality, including wastewater process imputation \cite{Mbamba2025EnhancingDQ}.

The promise of VAEs should be treated with some care. A VAE can learn a useful latent structure, but it can also produce overly smooth reconstructions if the objective rewards average plausibility more than local accuracy. In a temperature series, a smooth curve is not always a correct curve. The model may preserve the broad trend while missing a short spike or a sharp transition. This is one reason the present study uses a residual design. Instead of asking the VAE to reconstruct the entire signal from scratch, a deterministic prior supplies a strong first estimate and the LSTM-VAE learns corrections around it.

\subsection{Generative adversarial imputation models}

Generative Adversarial Imputation Nets (GAIN) adapt the adversarial idea to missing-data imputation. A generator fills missing entries, while a discriminator attempts to identify which entries were observed and which were imputed \cite{Yoon2018GAINMD}. Later surveys and extensions show that GAN-based imputation has become a substantial line of work \cite{Shahbazian2023GenerativeAN}. For time-series and sensor data, GAN variants have also been combined with random forests, anomaly detection, and all-sensor imputation settings \cite{Zhao2024SGADGANSG,Xie2025AnEG}.

GAN-based imputation is appealing because it tries to match the distribution of completed data rather than only minimise a pointwise error. That can matter if the imputed series should look realistic, not just close on average. Yet adversarial training is not a free lunch. It can be unstable, hard to reproduce, and sensitive to the balance between generator and discriminator. Also, a realistic-looking imputation is not always the value that actually occurred. For sensor reconstruction, this distinction matters: the completed curve should be plausible, but it should also be accurate at the masked timestamps.

This thesis does not use a GAN, but the GAN literature helps frame the broader generative-imputation landscape. It shows the move away from simple point replacement toward models that learn distributions and conditional structure. At the same time, it reinforces the need for fair quantitative evaluation. A generative model can look impressive in a plot while still underperforming a careful interpolation or KNN baseline on masked-point MAE.

\subsection{Hybrid and attention-based approaches}

Recent imputation work increasingly combines model families. Some methods pair statistical or machine-learning estimates with neural refinement. Others add attention mechanisms, positional encodings, graph structure, transfer learning, uncertainty modules, or diffusion processes. This shift seems sensible because missing-data problems are rarely solved by one cue alone. Local continuity, seasonal repetition, spatial correlation, and learned latent structure may all carry useful information.

Attention-based and position-aware models are especially relevant when the location of a missing value matters. In industrial processes, position-encoding denoising autoencoders explicitly incorporate position information \cite{Ou2024MissingDataIW}. In water-quality reconstruction, Bi-LSTM with self-attention and transfer learning has been used for continuous missing blocks \cite{Lyu2022}. Diffusion-based imputation methods, such as DiffPuter, represent a newer direction in which missing values are reconstructed through a generative denoising process \cite{Zhang2024DiffPuterED}. These methods may offer rich reconstructions, but they often increase training cost and implementation complexity.

Hybrid approaches are attractive for the present work because the underlying temperature signal has strong low-frequency structure. A deterministic prior such as interpolation or seasonal KNN already captures much of the obvious shape. The neural model can then focus on residual errors rather than relearning the entire curve. This design is less dramatic than a fully generative story, but it may be more practical: first use the simple method that is likely to be right most of the time, then allow the sequence model to adjust where the prior appears insufficient.

\section{LSTM and VAE Models for Time-Series Reconstruction}

The proposed model sits at the intersection of LSTM sequence modelling and VAE-style latent reconstruction. The two components address different parts of the imputation problem. The LSTM handles temporal ordering within windows. The VAE component encourages the model to represent the window through a compressed latent distribution. In combination, an LSTM-VAE can reconstruct sequences while learning a lower-dimensional description of their temporal behaviour.

\subsection{LSTM architectures for temporal dependencies}

An LSTM uses gates to control what information is kept, updated, and exposed at each timestep. This design was developed to reduce the vanishing-gradient problems that limited earlier recurrent networks \cite{Hochreiter1997}. For imputation, the relevant point is straightforward: the model can carry information across multiple timesteps instead of treating each timestamp as independent. A morning temperature reading may depend on the previous several readings, not just the value immediately before it.

Window design is important. A short window can make training easier but may miss slower trends. A long window can capture broader context but increases computation and may dilute local detail. Bidirectional models can use both past and future context, which is helpful for offline imputation inside a sequence, but less suitable for live imputation where future values have not arrived. The START, MIDDLE, END, and RANDOM settings in this thesis make these context differences visible. MIDDLE gaps can use information from both sides. START and END gaps cannot.

LSTM-based methods in the literature often perform well when the signal has repeated temporal patterns, but they are not automatically superior. Some forecasting comparisons find that simpler time-series or machine-learning models can outperform LSTM under certain conditions \cite{RoseroMontalvo2023TimeSeriesFT,Wu2024}. That does not undermine LSTM imputation; it just keeps the claim honest. A sequence model should earn its place through comparison against strong baselines.

\subsection{Latent variable modelling with VAEs}

Latent variable models assume that observed data can be explained through lower-dimensional hidden structure. In a temperature sequence, that hidden structure might correspond loosely to daily rhythm, seasonal level, local trend, and noise. A VAE-style model learns a distribution over such latent codes and uses sampled or mean latent values to reconstruct the input. This can support reconstruction under uncertainty because missing values are inferred from a learned representation rather than copied from a fixed rule.

For imputation, the latent representation may help in two ways. First, it can smooth over small corruptions by relying on broader sequence structure. Second, it can reduce the burden on the model when many input points are missing: instead of making each point independent, the decoder reconstructs a coherent window from a shared latent code. In practice, however, VAEs can also blur local behaviour. The KL regularisation that makes the latent space well behaved may also encourage conservative reconstructions. This is not a fatal flaw, but it suggests why the reconstruction loss, residual penalty, and prior design matter.

VAE applications in sensor and process imputation support this direction. Okafor and Delaney report VAE-based imputation as part of an IoT sensor calibration study \cite{Okafor2021}, while Mbamba et al. use deep VAEs in wastewater process data quality work \cite{Mbamba2025EnhancingDQ}. These studies suggest that VAEs can be useful beyond image-like data. Still, they also point back to domain context: an imputer for environmental temperature should be judged on temporal reconstruction, not only on whether the architecture is fashionable.

\subsection{LSTM-VAE designs for imputation}

An LSTM-VAE usually combines an LSTM encoder, latent distribution layers, and an LSTM decoder. The encoder reads an input window and summarises it into latent parameters. The latent code is sampled during training and passed to the decoder, which reconstructs the sequence. For missing-data imputation, the input can include the filled sequence, a binary mask, calendar features, and sometimes a first-pass imputation. The output can be the reconstructed value itself or a correction to a prior estimate.

The residual version used in this thesis is motivated by a fairly practical observation: if interpolation or seasonal KNN already gives a good answer, the neural network should not be forced to replace it completely. Learning a residual correction makes the neural component more conservative. It can improve the prior where useful, but the residual penalty and validation gate discourage unnecessary changes. This is especially relevant for smooth temperature data, where the best imputation under RANDOM missingness may simply be interpolation.

The LSTM-VAE design also fits the evaluation goal. The thesis is not trying to show that a deep model can fill missing values in the abstract. It asks whether a hybrid residual LSTM-VAE can outperform statistical and machine-learning baselines under shared masks. The literature suggests that such a comparison is necessary. Deep models may have more expressive capacity, but expressive capacity is not the same thing as better imputation on a particular sensor stream.

\section{Evaluation Strategies in the Literature}

Evaluation is one of the hardest parts of imputation research. With genuine missing values, the truth is unknown. With artificial missing values, the truth is known, but the missingness may be too clean or too unlike real failures. This tension appears throughout the literature and is especially visible in time-series work \cite{Farjallah2025EvaluationOM,Toye2025BenchmarkingMD}. A careful study needs to state exactly how values are hidden, which points are scored, which baselines are used, and what counts as improvement.

\subsection{Error metrics for imputation}

Mean Absolute Error (MAE) and Root Mean Squared Error (RMSE) are common choices for imputation evaluation. MAE measures average absolute deviation and is easy to interpret in the unit of the original variable. RMSE penalises larger errors more strongly because errors are squared before averaging. Many sensor-imputation studies report one or both metrics \cite{Agbo2022MissingDI,Choi2021,Lee2025ComparisonOM}.

The two metrics tell slightly different stories. MAE is useful when the concern is typical error. RMSE is more sensitive to occasional large mistakes, which may matter if sharp deviations are operationally important. In a temperature series, a method that is usually close but badly misses a short extreme event may have a reasonable MAE and a poor RMSE. Neither metric is sufficient on its own. That is why visual inspection remains useful: a curve can have a low average error but still look physically odd around a gap.

Another detail is where the metric is computed. In imputation, scoring the whole completed series can be misleading because observed points are often copied unchanged. A method may appear excellent simply because most points were never missing. The more meaningful evaluation computes error only at the masked positions. This thesis follows that approach by reporting masked-point MAE and RMSE.

\subsection{Masking strategies and benchmark design}

Benchmark design shapes the result. Missingness rate, gap length, missingness mechanism, mask location, and random seed can all change the ranking of methods. A single random 10\% mask may favour interpolation and KNN. A 50\% long block may favour models that use seasonal or global structure. A boundary gap may expose methods that quietly rely on having neighbours on both sides. Reviews and benchmark papers increasingly stress the need to vary missingness settings instead of reporting one number \cite{Miao2023AnES,Toye2025BenchmarkingMD,Zahmatkesh2026SpatioTemporalMD}.

The START, MIDDLE, END, and RANDOM masks used in this thesis are a deliberately simple response to that issue. START tests the absence of previous context. END tests the absence of future context. MIDDLE tests a long internal block. RANDOM tests scattered missingness. These masks do not cover every real mechanism, but they make the evaluation more informative than a single random mask. They also match realistic sensor problems: a device may fail at startup, stop reporting near the end of collection, drop a long block during outage, or lose isolated packets.

There is also the question of synthetic versus real missingness. Artificial masking gives ground truth, but real missingness gives realism. Farjallah et al. discuss the difficulty of evaluating imputation when ground truth is absent \cite{Farjallah2025EvaluationOM}. In this thesis, synthetic masking is chosen for fairness and measurement clarity. The limitation is acknowledged: the results indicate performance under controlled missingness, not a guarantee under every future field failure.

\subsection{Fair comparison across methods}

Fair comparison means more than listing several methods in the same table. The methods need to receive the same observed data, the same masks, and the same scoring rule. Preprocessing should not leak hidden values into training. Hyperparameter choices should be reasonable but not tailored unfairly to one method. If a neural model uses calendar features, a machine-learning baseline should not be denied equivalent time information unless the study is explicitly testing representation differences.

The literature contains many warnings in this direction, sometimes implicitly. Time-series benchmarks show that strong simple and machine-learning baselines can be difficult to beat \cite{Toye2025BenchmarkingMD,Luo2022}. IoT-edge comparisons remind us that computation time and resource use may matter alongside error \cite{Erhan2021CriticalCO}. Explainability-focused work suggests that choosing an imputer can also affect downstream model interpretation, not only predictive accuracy \cite{Erez2024}. These perspectives make the same general point: imputation methods should be evaluated as parts of a pipeline.

For this thesis, fairness is handled through shared masks, masked-point metrics, separate statistical and machine-learning baselines, and imputation-only runtime measurement. The comparison includes weak baselines, such as mean and median, but does not stop there. It also includes time-aware and lag-based baselines that are more credible competitors for a smooth sensor series. This makes the proposed model's win rate more meaningful.

\section{Summary of Research Gaps}

\subsection{Limitations in existing studies}

Several gaps in the literature motivate the present work. First, many studies still rely heavily on randomly scattered missingness. Random masks are useful, but they do not fully represent boundary gaps or long contiguous outages. Sensor failures are often block-like, and the lack of context at the start or end of a sequence can change the problem quite sharply \cite{Ahmed2022LongGM,Wu2022DataIF,Toye2025BenchmarkingMD}.

Second, deep imputation methods are sometimes evaluated against baselines that are too weak. A model that beats mean imputation has not necessarily shown much. For smooth physical signals, linear interpolation, KNN regression, gradient boosting, or iterative regression can be strong. Without those comparisons, the added value of a neural model remains uncertain. Recent benchmark and survey work makes this point indirectly by showing that method rankings vary across datasets and settings \cite{Miao2023AnES,Toye2025BenchmarkingMD}.

Third, some deep models reconstruct the whole signal even when a simple prior is already close. This can be wasteful and occasionally counterproductive. A residual approach may be better suited to sensor data where the baseline curve is often plausible but not perfect. The literature contains several hybrid directions, including combinations of statistical, machine-learning, attention, and generative components \cite{Lyu2022,Ou2024MissingDataIW,Zhang2024DiffPuterED}. What remains useful is a clear, controlled test of whether the hybrid idea actually improves masked-point accuracy for a specific IoT temperature series.

Fourth, runtime and practical cost are not always reported in a way that helps deployment thinking. A method may have good accuracy but require heavy training or slow inference. In edge settings, that matters \cite{Erhan2021CriticalCO}. This thesis does not claim to solve deployment, but it does separate imputation-only runtime from training so that the practical cost is at least visible.

\subsection{Motivation for the proposed approach}

The proposed hybrid residual LSTM-VAE is motivated by these gaps. It uses controlled START, MIDDLE, END, and RANDOM masks to avoid treating missingness as one uniform condition. It compares against both simple statistical baselines and stronger machine-learning baselines. It uses deterministic priors because interpolation and seasonal matching are not naive in this domain; they are serious competitors. The LSTM-VAE then learns residual corrections rather than replacing the prior wholesale.

This positioning is intentionally cautious. The thesis does not assume that deep learning is automatically superior. It asks a narrower and more useful question: when the same temperature series is masked in the same way, does the hybrid residual LSTM-VAE produce lower masked-point MAE and RMSE than methods a practitioner might actually use? If the answer is yes across gap types, the model has earned the extra complexity. If the advantage appears mainly for boundary or long-block gaps, that still tells us something practical. It suggests where sequence-based latent modelling may be worth the effort and where simpler imputers remain enough.

Overall, the literature supports three ideas that shape this study. Missingness in sensor time series is patterned and consequential. Simple and machine-learning baselines can be surprisingly competitive. Deep generative sequence models may help, but only when they are tested under fair, gap-aware evaluation. The next chapter builds the methodology around those ideas.
